% ============================================================================
%  00-map.tex — bản đồ bài học
%  Ngày tạo: 2026-09-09 | Sửa gần nhất: 2026-09-09 | Ôn gần nhất: chưa
% ============================================================================
\documentclass[hieu-suat.tex]{subfiles}
\begin{document}
\chapter*{Bản đồ bài học}
\addcontentsline{toc}{chapter}{Bản đồ bài học}
\markboth{Bản đồ bài học}{}

\begin{tomtat}
Bài học này trả lời một câu hỏi: \emph{người làm khoa học giỏi luyện tập cái gì, bằng cách nào, khi mà nghề của họ hầu như không cho phản hồi tử tế?} Nó được dựng từ 51 công trình gốc, sắp xếp lại theo mạch khái niệm chứ không theo thứ tự tôi đọc được. Xương sống của toàn bộ mười chương là một luận điểm duy nhất: mọi kỹ thuật luyện tập đã được chứng minh đều giả định một \textbf{vòng phản hồi nhanh, chính xác và lặp lại được}; nghiên cứu khoa học không có sẵn vòng đó, nên phần lớn công việc nâng hiệu suất chính là \emph{tự chế ra} nó. Nếu chỉ nhớ một điều: đừng đi tìm mẹo học nhanh hơn, hãy đi tìm chỗ nào trong công việc của mình đang không có phản hồi, rồi dựng phản hồi cho chỗ đó.
\end{tomtat}

\section*{A. Bài học này gồm những gì}

\begin{itemize}
\item \textbf{Chương 1 --- Nền tảng tối thiểu.} Cách đọc bằng chứng của tâm lý học nhận thức: cỡ hiệu ứng, tương quan so với can thiệp, phân tích tổng hợp, và vì sao một khuyến nghị nghe hợp lý vẫn có thể không có bằng chứng nào. Cộng bộ từ vựng dùng suốt các chương sau.
\item \textbf{Chương 2 --- Môi trường phản hồi.} Luyện tập có chủ đích thật ra đòi hỏi năm điều kiện; nghiên cứu vi phạm gần hết. Đây là chương gốc, mọi chương sau là một cách chữa một vi phạm cụ thể.
\item \textbf{Chương 3 --- Đọc và biểu diễn.} Chuyên gia khác người mới ở chỗ họ \emph{nhìn thấy} bài toán khác, không phải ở chỗ nhớ nhiều hơn. Cách đọc để xây biểu diễn theo cấu trúc sâu.
\item \textbf{Chương 4 --- Viết và viết lại.} Viết lại ở tầng nào thì mới là viết lại; quy trình viết nào có bằng chứng và quy trình nào chỉ là khẩu hiệu được lặp lại nhiều lần.
\item \textbf{Chương 5 --- Tự phản biện.} Vì sao tự tìm lỗi trong lập luận của mình khó về mặt cơ chế, và bốn kỹ thuật đã đo được là có tác dụng.
\item \textbf{Chương 6 --- Bài toán khó.} Kiểm soát và giám sát tiến trình: thứ phân biệt người giải được với người thông minh mà đi lạc ba tuần.
\item \textbf{Chương 7 --- Học và giữ.} Truy hồi, giãn cách, xen kẽ: xếp hạng theo bằng chứng, kèm những kỹ thuật phổ biến đã bị bác bỏ.
\item \textbf{Chương 8 --- Thực nghiệm trong CV/ML.} Chương đặc thù ngành: những cách một nhà nghiên cứu thị giác máy tính tự lừa mình bằng số, và quy trình chống lại.
\item \textbf{Chương 9 --- Chọn bài toán.} Số lượng, chất lượng, và việc chọn vấn đề quan trọng --- phần có bằng chứng yếu nhất nhưng hệ quả lớn nhất.
\item \textbf{Chương 10 --- Đo hiệu suất của chính mình.} Chỉ số tự đo, nền sinh học (ngủ, chú ý), và một chương trình mười hai tuần ghép mọi thứ lại.
\end{itemize}

\section*{B. Vì sao chia như vậy}

Cách chia hiển nhiên là theo \emph{kỹ năng}: đọc, viết, lập trình, thí nghiệm. Tôi không chia thế, vì cách đó giấu mất điều quan trọng nhất mà khảo sát tìm ra: các kỹ năng này không độc lập, chúng đều là những nỗ lực chữa cùng một khuyết tật của môi trường --- phản hồi đến quá muộn, quá nhiễu, hoặc không bao giờ đến \cite{hogarth2015}. Viết lại là cách lấy phản hồi từ chính bản thảo. Tự phản biện là cách lấy phản hồi khi chưa có ai để hỏi. Thiết kế thí nghiệm chặt là cách làm cho phản hồi từ dữ liệu trở nên đáng tin. Đọc là cách mượn phản hồi mà người khác đã trả giá để có.

Vì vậy Chương 2 là gốc, và các chương 3--8 là sáu vòng phản hồi khác nhau, xếp theo trình tự một dự án nghiên cứu thực sự đi qua: đọc để hiểu địa hình, viết để nghĩ, tự phản biện để không tự lừa, tấn công bài toán, giữ kiến thức đã học, rồi đo bằng thí nghiệm. Chương 9 và 10 lùi ra một bước: chọn cái gì để làm, và làm sao biết mình đang khá lên.

\section*{C. Phụ thuộc và thứ tự đọc}

\begin{figure}[H]\centering
\begin{tikzpicture}[node distance=0.55cm and 0.85cm,
  m/.style={draw=steel,fill=bluebg,rounded corners=2pt,inner sep=5pt,align=center,font=\small,text width=2.75cm},
  g/.style={draw=accent,fill=amberbg,rounded corners=2pt,inner sep=5pt,align=center,font=\small,text width=2.75cm},
  o/.style={draw=violetdark,fill=violetbg,rounded corners=2pt,inner sep=5pt,align=center,font=\small,text width=2.75cm},
  ar/.style={-{Latex[length=2.2mm]},draw=steel,thick},
  lb/.style={font=\scriptsize\itshape,color=reddark,align=center}]
\node[m] (c1) {\textbf{1.} Nền tảng\\đọc bằng chứng};
\node[g,right=1.4cm of c1] (c2) {\textbf{2.} Môi trường\\phản hồi};
\node[m,above right=0.2cm and 1.4cm of c2] (c3) {\textbf{3.} Đọc};
\node[m,below=of c3] (c4) {\textbf{4.} Viết lại};
\node[m,below=of c4] (c5) {\textbf{5.} Tự phản biện};
\node[m,right=1.1cm of c3] (c6) {\textbf{6.} Bài toán khó};
\node[m,right=1.1cm of c4] (c7) {\textbf{7.} Học và giữ};
\node[m,right=1.1cm of c5] (c8) {\textbf{8.} Thực nghiệm CV};
\node[o,below=1.0cm of c2] (c9) {\textbf{9.} Chọn bài toán};
\node[o,below=1.0cm of c9] (c10) {\textbf{10.} Đo bản thân};
\draw[ar] (c1) -- node[lb,above] {cỡ hiệu ứng} (c2);
\draw[ar] (c2) -- (c3); \draw[ar] (c2) -- (c4); \draw[ar] (c2) -- (c5);
\draw[ar] (c3) -- (c6); \draw[ar] (c4) -- (c7); \draw[ar] (c5) -- (c8);
\draw[ar] (c2) -- (c9); \draw[ar] (c9) -- (c10);
\draw[ar,dashed] (c5) to[bend right=12] (c4);
\end{tikzpicture}
\caption{Quan hệ phụ thuộc. Chương 2 (nền vàng) là gốc lập luận; các chương 3--8 đọc được theo thứ tự bất kỳ sau khi đã đọc chương 2. Mũi tên đứt: tự phản biện và viết lại nuôi nhau, nên đọc gần nhau.}
\end{figure}

\section*{D. Cốt lõi so với tham khảo}

\textbf{Cốt lõi} \muc{3} --- chương 2, 5, 7, 8. Đây là bốn chương có bằng chứng mạnh nhất \emph{và} đổi được cách làm việc hằng ngày. Đọc kỹ, quay lại nhiều lần.

\textbf{Nên nắm} \muc{2} --- chương 3, 4, 6, 10. Bằng chứng chắc, nhưng phần áp dụng cần tự điều chỉnh nhiều theo hoàn cảnh.

\textbf{Biết là có} \muc{1} --- chương 1 (đọc theo nhu cầu, quay lại khi gặp một con số không biết diễn giải) và chương 9 (bằng chứng yếu nhất trong cả bài học; đọc để có bản đồ, đừng coi là chỉ dẫn).

\section*{E. Quy ước trích dẫn}

Bài học tuân theo \code{../research-rules.md} mục 6.3. Mỗi ý mượn có trích dẫn ngay tại câu chứa nó, kèm định vị khi cần. Ba nhãn phân biệt nguồn của mỗi khẳng định và không bao giờ trộn trong một câu:

\begin{itemize}
\item \nD{} --- tác giả \textbf{đo được}: có số, có điều kiện thí nghiệm.
\item \nT{} --- tác giả \textbf{tuyên bố hoặc suy luận}, không đo trực tiếp.
\item \nM{} --- \textbf{tôi} suy ra, hoặc tôi ghép từ nhiều nguồn; kể cả khi ghép nghe hiển nhiên.
\end{itemize}

Câu không mang nhãn và không có trích dẫn là lời dẫn của tôi, không phải khẳng định thực nghiệm.

\end{document}
